\section{Observable classification and finite Poisson allocation}
\label{sec:uf-foundations}

This section proves the specialized upper-bound tools used below. The
ordinary classifier is the construction in the supplied manuscript
\cite{original,repo}; the weighted reference bank and multilevel
allocation are restated with their proofs so that the present paper
does not depend on the preceding unpublished notes. All likelihoods
refer to the actual observed binary prefix. A reference distribution
is used to evaluate that prefix, and never reveals or redraw the
latent state of a unit.

\subsection{A capped classifier that does not know the gap}

For a dyadic integer $M\ge1$, inspect prefix lengths
$1,2,4,\ldots,M$. At a prefix $n<M$, stop if
\[
 \left|\widehat p_n-\frac12\right|>
 t_{n,M}:=\sqrt{\frac{1+\log(M/n)}{n}},
\]
and output the side of $1/2$ containing $\widehat p_n$. At $M$, output
$\1_{\{\widehat p_M>1/2\}}$, with a negative output on ties.
Denote the output and work by $(Z_M^{\rm ord},T_M^{\rm ord})$.

\begin{lemma}[Ordinary conditional error and work]\label{lem:uf-ordinary}
For every $0<d\le1/2$ and either latent sign,
\begin{align}
 \Pp_{d,s}(Z_M^{\rm ord}\ne s)
 &\le \min\{1,(Md^2)^{-2}\},\label{eq:uf-ordinary-error}\\
 \E_{d,s}T_M^{\rm ord}
 &\le o_M(d):=\min\{M,17d^{-2}[1+\log_+(Md^2)]\}.
 \label{eq:uf-ordinary-work}
\end{align}
\end{lemma}

\begin{proof}
For independent Bernoulli observations of mean $p$, Hoeffding's
inequality gives
$\Pp(\widehat p_n-p\ge t)\le e^{-2nt^2}$ and the analogous lower-tail
bound \cite{hoeffding}. A wrong early output at a fixed inspected prefix
requires a deviation of at least $d+t_{n,M}$. Applying a union bound on
the underlying full tape, without conditioning on the stopping event,
therefore gives
\[
 \Pp_{d,s}(Z_M^{\rm ord}\ne s)
 \le e^{-2Md^2}
       +e^{-2}\sum_{\substack{n<M\\n\ {\rm dyadic}}}
                    (n/M)^2e^{-2nd^2}.
\]
For any dyadic sequence of positive $u$ values, the intervals
$[u/2,u]$ have disjoint interiors and
\[
 \int_{u/2}^{u}y e^{-2y}\,dy
 \ge\frac38u^2e^{-2u}.
\]
Consequently $\sum_u u^2e^{-2u}\le2/3$.
Writing $x=Md^2\ge1$ and using $x^2e^{-2x}\le e^{-2}$, the preceding
error bound is at most $(5/3)e^{-2}x^{-2}<x^{-2}$. For $x<1$, the
probability bound by one suffices.

Since the stopping time takes dyadic values,
\[
 \E_{d,s}T_M^{\rm ord}
 =1+\sum_{\substack{n<M\\n\ {\rm dyadic}}}
                      n\,\Pp_{d,s}(T_M^{\rm ord}>n).
\]
If $d>2t_{n,M}$, survival implies a deviation from the true mean of
at least $d/2$. Thus
\[
 \E_{d,s}T_M^{\rm ord}
 \le1+\sum_{\substack{n<M\\n\ {\rm dyadic}}}
       n\left[\1_{\{d\le2t_{n,M}\}}+2e^{-nd^2/2}\right].
\]
For $x=Md^2\ge1$, the first indicator implies
$nd^2\le4(1+\log x)$. Indeed, this is immediate if $nd^2<1$; otherwise
$\log(M/n)=\log(x/(nd^2))\le\log x$.
The corresponding dyadic sum is at most
$8d^{-2}(1+\log x)$.
Also
\[
 \int_{u/2}^{u}e^{-y/2}\,dy\ge\frac12u e^{-u/2}
 \quad\Longrightarrow\quad
 \sum_u u e^{-u/2}\le4.
\]
The exponential part of the work bound is therefore at most $8d^{-2}$.
Since $1\le d^{-2}$, total expected work is at most
$17d^{-2}(1+\log x)$. Finally $T_M^{\rm ord}\le M$ deterministically;
when $x<1$, the displayed minimum equals $M$.
\end{proof}

For any finite measure $\xi$ on $(0,1/2]$, write
\[
 \eta_\xi(M)=\int\min\{1,(Md^2)^{-2}\}\,\xi(dd),
 \qquad
 O_\xi(M)=\int o_M(d)\,\xi(dd).
\]
Both are homogeneous in $\xi$. The first is nonincreasing in $M$ and
the second is nondecreasing. Pointwise comparison also gives
$\eta_\xi(4M)\ge\eta_\xi(M)/16$.

\begin{lemma}[A homogeneous dyadic kernel estimate]\label{lem:uf-kernel}
For every finite measure $\xi$ and integer $K\ge1$,
\[
 \sum_{k=1}^{K}\sqrt{\eta_\xi(4^k)O_\xi(4^k)}
 \le208\int\min\{2^K,d^{-1}\}\,\xi(dd).
\]
In particular,
\begin{equation}\label{eq:uf-kernel-infinite}
 \sum_{k\ge1}\sqrt{\eta_\xi(4^k)O_\xi(4^k)}
 \le208\int d^{-1}\,\xi(dd),
\end{equation}
with extended-value interpretation.
\end{lemma}

\begin{proof}
Use bands $(2^{-j-1},2^{-j}]$, $j\ge1$, with masses $w_j$, and put
$c_j=2^jw_j$. Direct substitution into the two kernels yields
\[
 \eta_\xi(4^k)\le2^{-k}
 \left[\sum_{j\ge k}c_j2^{k-j}
              +16\sum_{j<k}c_j2^{3(j-k)}\right]=2^{-k}A_k,
\]
\[
 O_\xi(4^k)\le2^k
 \left[\sum_{j\ge k}c_j2^{k-j}
        +136\sum_{j<k}c_j2^{j-k}(1+k-j)\right]=2^kD_k.
\]
For the second display on $j<k$, use
$d^{-2}\le4\,2^{2j}$ and
$1+\log(4^kd^2)\le2(1+k-j)$.
On $j\ge k$, the deterministic cap gives the first sum exactly as
an upper bound.

Apply $\sqrt{A_kD_k}\le(A_k+D_k)/2$ and sum over $k\le K$.
For a fixed $j\le K$, the three relevant geometric sums obey
\[
 \sum_{k\le j}2^{k-j}\le2,\qquad
 \sum_{r\ge1}2^{-3r}=\frac17,\qquad
 \sum_{r\ge1}2^{-r}(1+r)=3.
\]
The coefficient of $c_j$ is thus at most
$2+8/7+204<208$. For $j>K$, only the first kernels occur, with
coefficient at most $2\,2^{K-j}$. Hence the result is at most
\[
 208\sum_{j\le K}2^jw_j+2\,2^K\sum_{j>K}w_j.
\]
For $j\le K$, $\min\{2^K,d^{-1}\}\ge2^j$ on the band; for $j>K$
it equals $2^K$. This proves the finite bound. Increasing $K$ and
using monotone convergence gives \eqref{eq:uf-kernel-infinite}.
All rearrangements involve nonnegative terms, so normalization and
moment assumptions are unnecessary.
\end{proof}

\subsection{A weighted bank compared with the full null mixture}

For a binary prefix $x_{1:n}$, let $P_{p,n}(x_{1:n})$ be its probability
under iid $\operatorname{Ber}(p)$ observations, and let
$P_{{\rm f},n}=P_{1/2,n}$. Choose any nonempty finite or countable set
of canonical bands
\[
 I_j=(b_j/2,b_j],\qquad b_j=2^{-j-1},
\]
and positive weights $\pi_j$ summing to one on these bands.
Set $r_j^-=b_j/2$ and $r_j^+=b_j/4$, and define
\[
 L_{j,n}^-=\frac{P_{1/2-r_j^-,n}}{P_{{\rm f},n}},
 \qquad
 M_n=\frac{\sum_j\pi_jP_{1/2+r_j^+,n}}
           {P_{-,\mu,n}},
 \qquad
 P_{-,\mu,n}=\int P_{1/2-d,n}\,\mu(dd).
\]
For $0<\epsilon<1$, stop negative when any
$L_{j,n}^-\ge1/(\epsilon\pi_j)$, and stop positive when
$M_n\ge1/(\rho\epsilon)$. Give negative crossings priority.
A zero denominator triggers a positive stop. Denote the bank stopping
time by $T_{\rm bank}$; it may be infinite outside its selected bands.

\begin{lemma}[Full-mixture likelihood bank]\label{lem:uf-mixture-bank}
The probability of a wrong bank stop is at most $2\rho\epsilon$
uniformly over $q\in\mathcal Q_\rho$, and its false-positive probability
is at most $\rho\epsilon$. These statements concern the event that the
bank stops and do not assume termination at every magnitude.

For $d\in I_j$ in a selected band,
\begin{equation}\label{eq:uf-bank-negative}
 \E_{d,-}T_{\rm bank}
 \le\frac3{4d^2}
       [1+\log(1/\epsilon)+\log(1/\pi_j)].
\end{equation}
Suppose a finite number $V(d)>0$ satisfies
\[
 \beta_\mu(n):=\int e^{-2nx^2}\,\mu(dx)\le V(d)
 \qquad\text{for every integer }n\ge d^{-2}.
\]
Then
\begin{equation}\label{eq:uf-bank-positive}
 \E_{d,+}T_{\rm bank}
 \le\frac6{d^2}
 \left[1+\log(1/\epsilon)+\log(1/\pi_j)
                   +\log_+\frac{V(d)}{\rho}\right].
\end{equation}
In particular the bank has finite conditional work on every selected
band whenever such an envelope is supplied.
\end{lemma}

\begin{proof}
A nonnegative supermartingale $Y_n$ with $Y_0=1$ has
$\Pp(\sup_nY_n\ge a)\le1/a$: stopping at the first crossing or a fixed
horizon gives $1\ge\E Y_{T\wedge m}\ge a\,\Pp(T\le m)$, and increasing
the horizon proves the claim. This elementary maximal inequality also
works when the process is a supermartingale only on supported prefixes.

Under a positive coin of gap $d$, the next multiplier of $L_{j,n}^-$
has mean $1-4r_j^-d\le1$. Its negative-crossing probability is at most
$\epsilon\pi_j$. A countable union bound gives conditional
negative-crossing probability at most $\epsilon$, so its positive-sign
contribution under any admissible prior is at most $\rho\epsilon$.

Under the full all-negative mixture, $M_n$ is a nonnegative
supermartingale on supported prefixes. If a prefix $x$ has positive
null probability, conditional expectation sums only children $xu$
with positive null probability; their numerator probabilities sum to
at most the numerator at $x$. This proves the supermartingale
inequality even when $\mu$ has mass at $d=1/2$. A zero-denominator prefix
has null probability zero. Thus a positive crossing has null-mixture
probability at most $\rho\epsilon$. Its actual negative-sign
contribution is no larger, since it multiplies the conditional
crossing probabilities by $1-q(d)\le1$. The two error claims follow.

For negative work, fix a selected-band gap $d\in I_j$ and
$r=r_j^-\in[d/2,d]$, with $r\le1/4$. The matching log likelihood has
mean increment
\[
 g(d,r)=\frac12\log(1-4r^2)
          +d\log\frac{1+2r}{1-2r}
 \ge-\frac83r^2+4dr\ge\frac43d^2.
\]
Here $\log(1-x)\ge-x/(1-x)$ and
$\log((1+x)/(1-x))\ge2x$ establish the first lower bound; the
concave quadratic takes the value $4d^2/3$ at both endpoints
$r=d/2,d$. Every upward log increment is at most $\log(3/2)<1$.
For the first correct-reference crossing $\tau$, the stopped-sum
identity at $\tau\wedge m$ gives
\[
 g(d,r)\E(\tau\wedge m)
 =\E\log L^-_{j,\tau\wedge m}
 \le1+\log(1/(\epsilon\pi_j)).
\]
Monotone convergence proves integrability of $\tau$ and the stated
work bound, since $T_{\rm bank}\le\tau$. No unbounded stopping identity
is invoked before integrability is established.

For positive work, let
$E_n=\{\widehat p_n-1/2\ge d/2\}$. On this event, for $x<1/2$,
\[
 \log\frac{P_{1/2-x,n}}{P_{{\rm f},n}}
 =\frac n2\log(1-4x^2)
   -n(\widehat p_n-\tfrac12)
                 \log\frac{1+2x}{1-2x}
 \le-2nx^2.
\]
For $x=1/2$ the likelihood is zero, since $E_n$ entails at least one
head. Therefore $P_{-,\mu,n}/P_{{\rm f},n}\le\beta_\mu(n)$.
The matching positive reference $r=r_j^+\in[d/4,d/2]$ satisfies,
on the same event,
\[
 \log\frac{P_{1/2+r,n}}{P_{{\rm f},n}}
 \ge n\left[-\frac83r^2+2dr\right]\ge nd^2/3.
\]
Set
\[
 U=1+\log(1/\epsilon)+\log(1/\pi_j)
                  +\log_+(V(d)/\rho),\qquad
 n_0=\left\lceil3d^{-2}U\right\rceil.
\]
For every $n\ge n_0$, $E_n$ implies
$M_n\ge\pi_je^{nd^2/3}/V(d)\ge1/(\rho\epsilon)$.
Thus $T_{\rm bank}>n$ implies $E_n^c$, whose positive-coin probability
is at most $e^{-nd^2/2}$. Summing the survival probabilities yields
\[
 \E_{d,+}T_{\rm bank}
 \le3d^{-2}U+2+2d^{-2}\le6d^{-2}U.
\]
The last inequality uses $U\ge1$ and $d^2\le1/4$.

For a countable bank, the positive mixture is a consistent probability
law and the negative stopping event is a measurable countable union.
At any fixed time only finitely many negative references can cross:
$L^-_{j,n}\le(3/2)^n$ requires
$\pi_j\ge1/[\epsilon(3/2)^n]$ for a crossing. All preceding arguments
therefore apply to finite and countable banks alike.
\end{proof}

\subsection{Racing procedures and retaining a genuine unit}

\begin{lemma}[Observable race]\label{lem:uf-race}
Suppose two classifiers run on the odd and even subsequences of one
unit's observation tape, with independent auxiliary seeds. Stop the
physical procedure when either classifier stops and use its output.
On completed counterfactual subsequences its work satisfies
\[
 T_{\rm race}\le2\min\{T_1,T_2\}.
\]
Its error event is contained in the union of the component wrong-stop
events. If either component has a finite deterministic cap, the race
terminates with a finite deterministic cap.
\end{lemma}

\begin{proof}
By physical time $2n$, each component has received $n$ observations.
If either component would stop after $n$ of its own observations, the
race stops no later than $2n$. This proves the pathwise work bound,
including an infinite stopping time for the other component.
The selected output is wrong only if the component that produces it
makes a wrong stop. Conditional on the latent $(D,S)$, both
subsequences are iid Bernoulli tapes of the correct mean, so all
component error and work bounds apply. The policy uses only the
physical observed prefix and its seeds.
\end{proof}

\subsection{Finite multilevel allocation at every positive budget}

\begin{lemma}[Independent Poisson allocation]\label{lem:uf-poisson}
Let $Z_0,\ldots,Z_J\in\{0,1\}$ be observable single-unit classifiers,
each with finite deterministic cap. Suppose uniformly over
$q\in\mathcal Q_\rho$ that their error probabilities are at most
$e_\ell$, their false-positive probabilities are at most $f_\ell$,
and their expected work is at most $C_\ell>0$.
Put
\[
 v_0=\rho+f_0,\qquad
 v_\ell=e_\ell+e_{\ell-1}\quad(\ell\ge1),\qquad
 k_0=C_0,\qquad k_\ell=C_\ell+C_{\ell-1}\quad(\ell\ge1).
\]
Assume $v_\ell>0$ and set $A_J=\sum_{\ell=0}^J\sqrt{v_\ell k_\ell}$.
For every $B>0$ there is an observable population estimator of
uniform expected work at most $B$ and worst RMS error at most
\begin{equation}\label{eq:uf-poisson-risk}
 \sqrt{e_J^2+A_J^2/B}.
\end{equation}
\end{lemma}

\begin{proof}
A level-zero correction uses a fresh genuine unit and returns
$\Delta_0=Z_0$. For $\ell\ge1$, a correction draws one fresh latent unit,
runs $Z_\ell$ on its observations, then runs $Z_{\ell-1}$ with a fresh
seed on the next unused observations of the same unit, and returns
$\Delta_\ell=Z_\ell-Z_{\ell-1}$.

Conditional on the latent state, the unused suffix after a finite
stopping time is an iid tape with the same Bernoulli mean: for each
possible stopped prefix, independence of its subsequent coordinates
gives the claim, and summing over the stopped prefixes preserves it.
This conditional argument is applied to successive unused blocks of
the same latent unit, even though the first block has random length.
The second classifier therefore has its original marginal distribution
under the same latent state. No new magnitude or sign is drawn between
the two classifiers. Their disagreement implies that at least one
misclassifies that common sign, and hence
\[
 \E_q\Delta_\ell^2\le v_\ell,\qquad
 \E_q W_\ell\le k_\ell,\qquad
 \sum_{\ell=0}^J\E_q\Delta_\ell=\E_qZ_J.
\]
At level zero use
$\Pp_q(Z_0=1)\le\Pp_q(S=1)+f_0\le\rho+f_0$.

Independently across levels, draw Poisson replication counts with means
\[
 n_\ell=\frac B{A_J}\sqrt{\frac{v_\ell}{k_\ell}}.
\]
For each replication use an independent genuine unit, and return
\[
 \widehat\theta
 =\sum_{\ell=0}^J\frac1{n_\ell}
       \sum_{i=1}^{N_\ell}\Delta_{\ell,i},
 \qquad N_\ell\sim\operatorname{Poisson}(n_\ell).
\]
Its mean is $\E_qZ_J$, so its absolute bias is at most $e_J$.
The expected work is at most
$\sum_\ell n_\ell k_\ell=B$.

For independent marks $\Delta$ and a Poisson count of mean $n$,
conditioning on the count gives
$\Var(\sum_{i=1}^{N}\Delta_i)=n\E\Delta^2$.
Consequently the independent level variances sum to at most
\[
 \sum_{\ell=0}^J\frac{v_\ell}{n_\ell}=\frac{A_J^2}{B}.
\]
Bias squared plus variance proves \eqref{eq:uf-poisson-risk}.
All counts are almost surely finite, all component caps are finite,
and all first and subsequent observations are included in $W_\ell$.
\end{proof}

The allocation may use Poisson means below one; this is valid for
every $B>0$. The bound is uniform over the original rarity class.
Finite-level bounds can be optimized by taking an infimum over
finite constructions without asserting the existence of a single
infinite-level algorithm.

